\section{Discussion}\label{sec:discussion}
% DRAFT for the author to rewrite. Every statement points to where it is shown in the paper.

For modular addition, the number of features a maximum-margin network uses is set by the exponent $s=\nu/b$: every frequency at $b=\nu$, and $\Theta(\log p)$ frequencies for every $b<\nu$ (Theorems~\ref{thm:cyclic} and~\ref{thm:count}). The same programme recovers the dense maxima of \citet{morwani2024feature} and \citet{li2025nsum} at $b=\nu$, and in trained networks the count moves with $s$ in the direction the programme predicts (Section~\ref{sec:experiments}).

\paragraph{Limitations.} The exact optimum of Theorem~\ref{thm:cyclic} is proved for interference-free networks. For all networks we prove only the universal bound and the order of the optimum (Theorems~\ref{thm:universal} and~\ref{thm:scaling}); whether interference can raise the margin for $2\le b<\nu$ is open (Appendix~\ref{app:conjecture}). The results concern maximum-margin solutions, which regularized cross-entropy reaches only at its global minimizers as $\lambda\to0$ \citep{wei2019regularization}; gradient descent need not find them. Accordingly, trained networks carry some interference and are not maximizers of the programme, so our experiments test trends in the count rather than exact values: they do not resolve the finer plateaus of $\Kstar$ or its non-monotone steps in $p$, and they fall further below $\Kstar$ as the degree grows. For power networks we prove the same upper bounds (with constant $n!B_p$), but they are known to be attained only at $b=\nu$ \citep{li2025nsum}, so the exact optimum and the logarithmic laws are proved for product networks only (Appendix~\ref{app:conjecture}). Finally, the theory covers sums of such neurons over the cyclic group $\sZ_p$; networks with a shared linear layer in front of the neurons are priced by a nuclear norm and fall outside it (Appendix~\ref{app:models}), and other groups are not treated here.
